\Needspace{7\baselineskip}
\subsection{Row 1: Refinement map}
\label{row:1}
\textbf{Classification:} Exact.
\par\smallskip\textbf{Declaration targets}\par
{\small\noindent\code{NR.children}\par}
{\small\noindent\code{NR.mem_children}\par}
{\small\noindent\code{NR.Refines}\par}
\smallskip\textbf{Checked scope.} For finite complete projector families over a semiring module, \code{Refines} records a surjective parent map and equality of every coarse projector with the finite sum over its fine children.
\par\smallskip\textbf{Boundary.} No refinement is constructed from arbitrary projector families; no infinite or topological sum is asserted.
\medskip
\Needspace{7\baselineskip}
\subsection{Row 2: Nested refinement theorem}
\label{row:2}
\textbf{Classification:} Stronger.
\par\smallskip\textbf{Declaration targets}\par
{\small\noindent\code{NR.finitePacket_decomposition_unique}\par}
{\small\noindent\code{NR.finitePacket_decomposition_ext}\par}
{\small\noindent\code{NR.Refines.coarse_comp_fine_eq_self}\par}
{\small\noindent\code{NR.Refines.coarse_comp_fine_eq_zero}\par}
{\small\noindent\code{NR.Refines.fine_comp_coarse_eq_self}\par}
{\small\noindent\code{NR.Refines.fine_comp_coarse_eq_zero}\par}
{\small\noindent\code{NR.Refines.coarse_packet_reconstruction}\par}
{\small\noindent\code{NR.Refines.child_decomposition_unique}\par}
{\small\noindent\code{NR.Refines.trans}\par}
\smallskip\textbf{Checked scope.} Parent/child absorption holds in both composition orders and refinements compose. More generally, any finite sum of vectors lying in mutually orthogonal packet ranges has uniquely determined components, without a refinement or completeness hypothesis.
\par\smallskip\textbf{Boundary.} Uniqueness is only asserted on the selected finite set and requires each component to lie in its stated projector range. All index types and sums are finite.
\medskip
\Needspace{7\baselineskip}
\subsection{Row 3: Exact coordinate refinement}
\label{row:3}
\textbf{Classification:} Stronger.
\par\smallskip\textbf{Declaration targets}\par
{\small\noindent\code{NR.refineCoordinates}\par}
{\small\noindent\code{NR.coarsenCoordinates}\par}
{\small\noindent\code{NR.refineCoordinates_self}\par}
{\small\noindent\code{NR.refineCoordinates_comp}\par}
{\small\noindent\code{NR.coordinateRefinementEquiv_trans}\par}
{\small\noindent\code{NR.Refines.refineCoordinates_component}\par}
{\small\noindent\code{NR.Refines.coarsenCoordinates_component}\par}
\smallskip\textbf{Checked scope.} Analysis after synthesis defines cross-level linear maps for any two finite complete projector families; these maps obey identity and cocycle laws through every intermediate family. A \code{Refines} witness gives the parent/child component formulas.
\par\smallskip\textbf{Boundary.} These are coherent changes of exact coordinates, not truncation or dimensional reduction operators; the cocycle uses exact completeness at each finite level.
\medskip
\Needspace{7\baselineskip}
\subsection{Row 4: Perfect reconstruction across levels}
\label{row:4}
\textbf{Classification:} Stronger.
\par\smallskip\textbf{Declaration targets}\par
{\small\noindent\code{NR.coarsenCoordinates_refineCoordinates}\par}
{\small\noindent\code{NR.refineCoordinates_coarsenCoordinates}\par}
{\small\noindent\code{NR.coordinateRefinementEquiv}\par}
{\small\noindent\code{NR.coordinateOperator}\par}
{\small\noindent\code{NR.coordinateOperator_comp}\par}
{\small\noindent\code{NR.coordinateOperator_naturality}\par}
\smallskip\textbf{Checked scope.} Cross-level refinement and coarsening are mutual inverses over a semiring module. The typed coordinate representation of every linear map preserves identities and composition and commutes with independent exact coordinate changes on source and target.
\par\smallskip\textbf{Boundary.} The inverse laws need completeness of the projector families but no parent map. This is finite algebraic naturality, not convergence for an infinite packet tower.
\medskip
\Needspace{7\baselineskip}
\subsection{Row 5: Transfer-block refinement}
\label{row:5}
\textbf{Classification:} Stronger.
\par\smallskip\textbf{Declaration targets}\par
{\small\noindent\code{NR.transferBlock_refinement_between}\par}
{\small\noindent\code{NR.transferBlock_refinement}\par}
{\small\noindent\code{NR.transferBlock_refinement_and_reconstruction}\par}
\smallskip\textbf{Checked scope.} Source and target projector families may live on different modules and refine independently; every coarse linear-map block is the double finite sum of its fine descendant blocks. The direct traceability wrapper also records reconstruction of the same map from all coarse blocks and from all fine blocks. The endomorphism formula is a specialization.
\par\smallskip\textbf{Boundary.} No sparse extraction algorithm, infinite block matrix, or norm convergence is proved here.
\medskip
\Needspace{7\baselineskip}
\subsection{Row 6: Residue projector algebra}
\label{row:6}
\textbf{Classification:} Exact.
\par\smallskip\textbf{Declaration targets}\par
{\small\noindent\code{R.residueProj_idempotent}\par}
{\small\noindent\code{R.residueProj_orthogonal}\par}
{\small\noindent\code{R.residueProj_complete}\par}
\smallskip\textbf{Checked scope.} For every positive modulus, coefficient residue selectors on \code{Nat -> A} are idempotent, mutually annihilating, and complete over an arbitrary semiring and module.
\par\smallskip\textbf{Boundary.} This is the division-free coefficient selector, not a character-average formula.
\medskip
\Needspace{7\baselineskip}
\subsection{Row 7: Root-of-unity realization}
\label{row:7}
\textbf{Classification:} Stronger / certificate-explicit.
\par\smallskip\textbf{Declaration targets}\par
{\small\noindent\code{R.ringFourierResidueKernel}\par}
{\small\noindent\code{R.ringFourierKernel_eq_residueKernel}\par}
{\small\noindent\code{R.FiniteRingFourierOrthogonalityCertificate}\par}
{\small\noindent\code{R.FiniteRingFourierOrthogonalityCertificate.toRingCertificate}\par}
{\small\noindent\code{R.ringFourierAverage}\par}
{\small\noindent\code{R.ringFourierAverage_eq_residueProj_of_finite_certificate}\par}
{\small\noindent\code{R.primitiveRoot_fourierOrthogonalityCertificate}\par}
{\small\noindent\code{R.fourierAverage_eq_residueProj_of_primitiveRoot}\par}
\smallskip\textbf{Checked scope.} Over any commutative semiring, \code{omega^N = 1} reduces the arbitrary coefficient address to a finite residue address. Consequently an explicit finite \code{N x N} orthogonality table identifies the division-free Fourier average with the residue selector. Over a field, a positive-order primitive root automatically supplies the normalized inverse-root certificate.
\par\smallskip\textbf{Boundary.} The general semiring branch still requires the finite orthogonality table; root-of-unity order alone is not enough in rings with zero divisors. The primitive-root construction is field-only.
\medskip
\Needspace{7\baselineskip}
\subsection{Row 8: Polyphase refinement}
\label{row:8}
\textbf{Classification:} Exact.
\par\smallskip\textbf{Declaration targets}\par
{\small\noindent\code{R.polyphase_refinement}\par}
\smallskip\textbf{Checked scope.} For positive \code{N} and \code{k}, the modulus-\code{N} selector is the finite sum of exactly those modulus-\code{N*k} selectors that merge to it, over any semiring module.
\par\smallskip\textbf{Boundary.} The theorem concerns exact coefficient masks; it gives no approximation or discarded-packet estimate.
\medskip
\Needspace{7\baselineskip}
\subsection{Row 9: Dyadic-odd packet factorization}
\label{row:9}
\textbf{Classification:} Exact.
\par\smallskip\textbf{Declaration targets}\par
{\small\noindent\code{R.crtEquiv}\par}
{\small\noindent\code{R.crtComponents_bijective}\par}
{\small\noindent\code{R.residueProj_crt}\par}
{\small\noindent\code{R.residueProj_crt_comm}\par}
{\small\noindent\code{R.dyadicOddEquiv}\par}
{\small\noindent\code{R.dyadicOddComponents_bijective}\par}
{\small\noindent\code{R.residueProj_dyadicOdd}\par}
\smallskip\textbf{Checked scope.} Positive coprime moduli give a bijection of residue labels and the composite selector in both composition orders. For positive odd \code{M}, this specializes to \code{2^a * M}.
\par\smallskip\textbf{Boundary.} This is coordinate CRT factorization. It does not identify the denominator or lifting behavior of spectral projectors.
\medskip
\Needspace{7\baselineskip}
\subsection{Row 10: Certified weighted inverse}
\label{row:10}
\textbf{Classification:} Stronger / certificate-explicit.
\par\smallskip\textbf{Declaration targets}\par
{\small\noindent\code{WT.TwoSidedInverse}\par}
{\small\noindent\code{WT.finiteGeometricInverse_spec_of_pow_eq_zero}\par}
{\small\noindent\code{WT.diagonalTransferInverse_spec}\par}
{\small\noindent\code{WT.analyticGeometricInverse_spec}\par}
{\small\noindent\code{WT.analyticDiagonalTransferInverse_spec}\par}
\smallskip\textbf{Checked scope.} In a possibly noncommutative ring, an explicit two-sided inverse for \code{D} and nilpotence of the normalized perturbation \code{Dinv*H} give a finite two-sided inverse of \code{D+H}. In a complete normed ring, \code{norm (Dinv*H) < 1} gives the separate infinite Neumann inverse.
\par\smallskip\textbf{Boundary.} The manuscript perturbation is \code{H}; the Lean declaration's parameter name \code{K} has no fraction-field meaning. There is no unrestricted convergence theorem for an abstract complete filtered ring. The analytic candidate is noncomputable.
\medskip
\Needspace{7\baselineskip}
\subsection{Row 11: Finite precision and weighted pruning}
\label{row:11}
\textbf{Classification:} Stronger / certificate-explicit.
\par\smallskip\textbf{Declaration targets}\par
{\small\noindent\code{WT.MultiplicativeFiltration}\par}
{\small\noindent\code{WT.PrecisionCertificate}\par}
{\small\noindent\code{WT.retainedFixedWordSum}\par}
{\small\noindent\code{WT.prunedWordInverse}\par}
{\small\noindent\code{WT.map_finiteGeometricInverse_eq_prunedWordInverse}\par}
{\small\noindent\code{WT.minWeight_mul_length_le_wordWeight}\par}
{\small\noindent\code{WT.map_sum_pow_eq_zero_at_scaled_cutoff}\par}
{\small\noindent\code{WT.prunedWordInverse_spec_of_minWeight}\par}
{\small\noindent\code{WT.finitePrecisionPrunedInverse_of_weightedEdges}\par}
{\small\noindent\code{WT.finitePrecisionPrunedDiagonalTransferInverse_of_minWeight}\par}
{\small\noindent\code{WT.finitePrecisionPrunedDiagonalTransferInverse_of_weightedEdges}\par}
\smallskip\textbf{Checked scope.} The mapped finite geometric inverse equals an explicit noncommutative word sum retaining exactly the words of cumulative weight below \code{L}. If every edge weight is at least \code{m} and \code{L <= m*N}, the \code{N}th mapped power vanishes and the pruned word sum is a two-sided inverse; the corresponding diagonal-plus-transfer formula follows from supplied \code{Dinv} and \code{Dinv*H = sum edge}.
\par\smallskip\textbf{Boundary.} The filtration, precision quotient map, edge-membership data, diagonal inverse, and edge expansion remain inputs. No DVR quotient, quotient map, basis-level modular matrix routine, or extracted pruning implementation is constructed.
\medskip
\Needspace{7\baselineskip}
\subsection{Row 12: Finite-graph transfer-zeta rationality}
\label{row:12}
\textbf{Classification:} Automatic.
\par\smallskip\textbf{Declaration targets}\par
{\small\noindent\code{TZ.coefficientAlgorithm_correct}\par}
{\small\noindent\code{TZ.coefficientAlgorithm_firstEdge}\par}
{\small\noindent\code{TZ.weightedAdjacency}\par}
{\small\noindent\code{TZ.transferMatrix}\par}
{\small\noindent\code{TZ.transferDenominator}\par}
{\small\noindent\code{TZ.transferNumerator}\par}
{\small\noindent\code{TZ.transferDenominator_coeff_zero}\par}
{\small\noindent\code{TZ.pathZeta_eq_matrixPathZeta}\par}
{\small\noindent\code{TZ.transferDenominator_mul_pathZeta}\par}
{\small\noindent\code{TZ.finite_graph_pathZeta_rationality}\par}
{\small\noindent\code{TZ.finite_graph_transferZeta_trace}\par}
{\small\noindent\code{TZC.certificate_explicit_rationality}\par}
\smallskip\textbf{Checked scope.} For every finite positive-weight directed multigraph and source/target pair, finite walk enumeration satisfies the first-edge recurrence. Lean constructs \code{Q = det(I-A(X))} and \code{P} as the corresponding adjugate entry, proves \code{Q(0)=1}, identifies the independently enumerated \code{pathZeta} with the matrix inverse entry, and proves \code{Q*Z=P} without a rationality certificate. A direct wrapper packages the manuscript's enumeration, recurrence, normalization, nonvanishing, and rational identity in one declaration. The older coefficient-certificate route remains available.
\par\smallskip\textbf{Boundary.} Walks may repeat vertices and edges; no simple-path restriction is imposed. Parallel edges are distinct edge values. Determinant and adjugate are noncomputable Lean definitions here; no denominator cancellation, minimality, extracted symbolic implementation, or analytic convergence claim is made.
\medskip
\Needspace{7\baselineskip}
\subsection{Row 13: Linear recurrence}
\label{row:13}
\textbf{Classification:} Automatic.
\par\smallskip\textbf{Declaration targets}\par
{\small\noindent\code{TZ.finite_graph_constant_coefficient_recurrence}\par}
{\small\noindent\code{TZC.constant_coefficient_recurrence}\par}
\smallskip\textbf{Checked scope.} The automatically constructed determinant denominator gives an antidiagonal constant-coefficient recurrence for every \code{n} beyond the degree of the canonical adjugate numerator. A separately supplied \code{P,Q} certificate yields the same form through the legacy route.
\par\smallskip\textbf{Boundary.} The canonical recurrence is not claimed minimal, and no recurrence is guessed or certified from finitely many sampled coefficients alone.
\medskip
\Needspace{7\baselineskip}
\subsection{Row 14: Evaluated Bezout spectral projector}
\label{row:14}
\textbf{Classification:} Stronger / certificate-explicit.
\par\smallskip\textbf{Declaration targets}\par
{\small\noindent\code{RS.EvaluatedBezoutCertificate}\par}
{\small\noindent\code{RS.separatorNumerator_eq_scalar_sub_of_bezout}\par}
{\small\noindent\code{RS.separatorNumerator_mul_self_of_bezout}\par}
{\small\noindent\code{RS.separator_idempotent_of_bezout}\par}
{\small\noindent\code{RS.separatorNumerator_apply_of_mem_ker_F_of_bezout}\par}
{\small\noindent\code{RS.separator_apply_of_mem_ker_F_of_bezout}\par}
{\small\noindent\code{RS.separator_apply_of_mem_ker_G}\par}
{\small\noindent\code{RS.scalar_mul_separator_of_ne_zero}\par}
{\small\noindent\code{RS.scalar_mul_separator_preservesLattice_of_ne_zero}\par}
\smallskip\textbf{Checked scope.} Over a field, the evaluated equality \code{U*F + V*G = r*1}, \code{F*G=0}, and only the two commutations actually used prove the denominator-free identity \code{(VG)^2=r(VG)}; adding \code{r!=0} gives idempotence of \code{r^-1*VG}. The kernel actions need only the displayed evaluated Bezout or kernel hypotheses, and lattice preservation of the numerator follows when \code{V,G} preserve the lattice.
\par\smallskip\textbf{Boundary.} The evaluated Bezout equality and annihilation are supplied. No polynomial Bezout identity, resultant, characteristic factorization, or Cayley-Hamilton instance is constructed; nilpotence alone is not used to infer exact kernel action.
\medskip
\Needspace{7\baselineskip}
\subsection{Row 15: Separation weight}
\label{row:15}
\textbf{Classification:} Stronger / certificate-explicit.
\par\smallskip\textbf{Declaration targets}\par
{\small\noindent\code{RS.ScalarPreservesLattice}\par}
{\small\noindent\code{RS.IsLatticeUnit}\par}
{\small\noindent\code{RS.IsSeparationWeight}\par}
{\small\noindent\code{RS.IsSeparationWeight.mono}\par}
{\small\noindent\code{RS.separationWeight}\par}
{\small\noindent\code{RS.separationWeight_clears}\par}
{\small\noindent\code{RS.separationWeight_le_of_clears}\par}
{\small\noindent\code{RS.pow_smul_separator_eq_unit_inv_smul_separatorNumerator}\par}
{\small\noindent\code{RS.separator_isSeparationWeight_of_unit_mul_pow_eq_of_inv_preserves}\par}
{\small\noindent\code{RS.separator_hasFiniteSeparationWeight_of_eq_unit_mul_pow}\par}
{\small\noindent\code{RS.separator_separationWeight_le_of_eq_unit_mul_pow}\par}
\smallskip\textbf{Checked scope.} \code{IsSeparationWeight pi E lattice s} means that \code{(pi^s)E} preserves the chosen lattice; weights are upward closed when scalar multiplication by \code{pi} preserves it. If \code{r=u*pi^s}, then \code{pi^s E=u^-1 VG}; preservation by the inverse unit and numerator proves clearing, while a bundled lattice-unit certificate gives existence and the least-weight upper bound \code{separationWeight<=s}.
\par\smallskip\textbf{Boundary.} The unit action on the chosen lattice is explicit data. The least weight is noncomputable via \code{Nat.find}; no valuation, polynomial resultant construction, search procedure, or claim that the displayed exponent is minimal is supplied.
\medskip
\Needspace{7\baselineskip}
\subsection{Row 16: Projector precision delay}
\label{row:16}
\textbf{Classification:} Exact from certificate.
\par\smallskip\textbf{Declaration targets}\par
{\small\noindent\code{RS.CongruentModLattice}\par}
{\small\noindent\code{RS.projector_precision_delay}\par}
\smallskip\textbf{Checked scope.} For an arbitrary field scalar \code{pi}, additive lattice, and clearing-weight certificate, an explicit congruence witness at depth \code{m+s} maps to one at depth \code{m}.
\par\smallskip\textbf{Boundary.} Congruence is witness-based scalar-lattice membership; no topology, completion, digit representation, or valuation API is required or concluded.
\medskip
\Needspace{7\baselineskip}
\subsection{Row 17: Odd-adic versus dyadic sign behavior}
\label{row:17}
\textbf{Classification:} Stronger / scoped realization.
\par\smallskip\textbf{Declaration targets}\par
{\small\noindent\code{SD.signInvolution_mul_self}\par}
{\small\noindent\code{SD.positiveSignProjector_idempotent}\par}
{\small\noindent\code{SD.negativeSignProjector_idempotent}\par}
{\small\noindent\code{SD.signProjectors_complete}\par}
{\small\noindent\code{SD.HasExactPLocalDenominatorWeight}\par}
{\small\noindent\code{SD.positiveSignProjector_exact_pLocal_weight_zero_of_prime_ne_two}\par}
{\small\noindent\code{SD.negativeSignProjector_exact_pLocal_weight_zero_of_prime_ne_two}\par}
{\small\noindent\code{SD.signFactors_resultant}\par}
{\small\noindent\code{SD.signProjectors_resultant_dyadic_weight_endpoint}\par}
{\small\noindent\code{SD.padicValNat_two_at_odd_prime}\par}
{\small\noindent\code{SD.oddAdicVersusDyadicSignBehavior_trace}\par}
\smallskip\textbf{Checked scope.} The displayed rational \code{2 x 2} involution and sign projectors are checked entrywise. Each projector has exact local denominator weight zero at every prime other than \code{2} and exact \code{2}-local weight one. Lean also computes \code{resultant(X-1,X+1)=2} and identifies its \code{2}-adic valuation with the exact local weight of both projectors. The direct trace wrapper combines both local regimes and the resultant statements.
\par\smallskip\textbf{Boundary.} This is an explicit rational-entry, natural-valuation model. Reduction modulo \code{2} applies to the characteristic factors, not to the rational projector matrices containing \code{1/2}. Lean does not construct \code{Q_p}, its valuation ring, or a general equivalence between rational \code{p}-integrality and preservation of an arbitrary completed lattice.
\medskip
\Needspace{7\baselineskip}
\subsection{Row 18: Concrete ZMod Smith-weighted congruence solve}
\label{row:18}
\textbf{Classification:} Automatic at concrete scope / certificate-explicit extension.
\par\smallskip\textbf{Declaration targets}\par
{\small\noindent\code{SD.smithZModMulHom}\par}
{\small\noindent\code{SD.SmithZModAdmissible}\par}
{\small\noindent\code{SD.smithZModFiber_nonempty_iff}\par}
{\small\noindent\code{SD.smithZModFiber_card_of_admissible}\par}
{\small\noindent\code{SD.smithZModCoordinateCertificate}\par}
{\small\noindent\code{SD.smithZMod_global_solvable_iff}\par}
{\small\noindent\code{SD.smithZMod_global_conditional_count}\par}
{\small\noindent\code{SD.SmithSolvabilityCertificate}\par}
{\small\noindent\code{SD.SmithCardinalityCertificate}\par}
{\small\noindent\code{SD.smith_conditional_product_count}\par}
\smallskip\textbf{Checked scope.} For \code{p>1}, Lean proves directly over \code{ZMod(p^depth)} that \code{p^weight*x=target} is solvable exactly when \code{p^min(weight,depth)} divides \code{target.val}, and every admissible fibre has cardinality \code{p^min(weight,depth)}. It then proves the coordinatewise criterion and the exact conditional count \code{p^(sum min(weight,depth))} for finite diagonal systems, including cardinality zero when a coordinate is inadmissible. The split certificate API gives the analogous product statement for abstract finite fibres, with the residue cardinality supplied as \code{q}.
\par\smallskip\textbf{Boundary.} The counted manuscript theorem is the concrete \code{ZMod(p^L)} result with \code{q=p}. The arbitrary-DVR, residue-size-\code{q} extension is classical/certificate-explicit, not derived by Lean. Diagonal weights and targets are inputs; Lean does not construct Smith normal form. A literal zero diagonal factor is represented at precision \code{L} by a weight at least \code{L}, or handled by an explicit abstract coordinate certificate.
\medskip
\Needspace{7\baselineskip}
\subsection{Row 19: Power-defect identity}
\label{row:19}
\textbf{Classification:} Stronger.
\par\smallskip\textbf{Declaration targets}\par
{\small\noindent\code{TLD.defect}\par}
{\small\noindent\code{TLD.power_defect_identity_all}\par}
{\small\noindent\code{TLD.power_defect_identity}\par}
{\small\noindent\code{TLD.powers_compatible_of_defect_eq_zero}\par}
{\small\noindent\code{TLD.polynomialAction}\par}
{\small\noindent\code{TLD.polynomialAction_compatible_of_defect_eq_zero}\par}
{\small\noindent\code{TLD.continuousPolynomialAction_compatible_of_defect_eq_zero}\par}
{\small\noindent\code{LD.power_defect_identity_all}\par}
\smallskip\textbf{Checked scope.} The primary theorem is typed for distinct fine and coarse semiring modules and proves the ordered telescoping identity for every natural \code{n}, including zero. Zero one-step defect gives exact compatibility with every power and every finite polynomial functional calculus; the continuous version is also proved.
\par\smallskip\textbf{Boundary.} This is linear algebra and finite polynomial calculus; no nonlinear higher-jet, operator-tree defect expansion, or analytic functional calculus is asserted.
\medskip
\Needspace{7\baselineskip}
\subsection{Row 20: Power-defect norm bound}
\label{row:20}
\textbf{Classification:} Stronger.
\par\smallskip\textbf{Declaration targets}\par
{\small\noindent\code{TLD.continuousDefect}\par}
{\small\noindent\code{TLD.continuous_power_defect_identity_all}\par}
{\small\noindent\code{TLD.continuous_end_pow_norm_le}\par}
{\small\noindent\code{TLD.continuous_norm_power_defect_le}\par}
{\small\noindent\code{TLD.continuous_norm_power_defect_le_geometric_majorant}\par}
{\small\noindent\code{LD.norm_power_defect_le_raw}\par}
{\small\noindent\code{LD.norm_power_defect_le}\par}
\smallskip\textbf{Checked scope.} Over a nontrivially normed field, continuous maps between distinct normed fine and coarse spaces satisfy the exact finite-sum operator-norm bound for every natural \code{n}, including zero. If both dynamics have norm at most \code{rho}, Lean derives the closed majorant \code{n*rho^(n-1)*norm(defect)}. No completeness assumption on the spaces is needed.
\par\smallskip\textbf{Boundary.} These are operator-norm estimates, not spectral-radius estimates, equality statements, or nonlinear defect bounds. The majorant takes \code{rho} and both norm inequalities as hypotheses.
\medskip
\Needspace{7\baselineskip}
\subsection{Row 21: Resolvent-defect identity}
\label{row:21}
\textbf{Classification:} Stronger / certificate-explicit.
\par\smallskip\textbf{Declaration targets}\par
{\small\noindent\code{TLD.TwoSidedInverse}\par}
{\small\noindent\code{TLD.scalarShift_resolvent_identity}\par}
{\small\noindent\code{TLD.resolvents_compatible_of_defect_eq_zero}\par}
{\small\noindent\code{TLD.continuousScalarShift}\par}
{\small\noindent\code{TLD.continuous_scalarShift_resolvent_identity}\par}
{\small\noindent\code{TLD.continuous_norm_resolvent_defect_le}\par}
{\small\noindent\code{TLD.continuous_resolvents_compatible_of_defect_eq_zero}\par}
{\small\noindent\code{LD.resolvent_defect_identity}\par}
\smallskip\textbf{Checked scope.} Algebraically, supplied two-sided inverse witnesses give the typed resolvent identity on distinct modules. For continuous maps, the proof needs only \code{(zI-Tf)Rf=I} and \code{Rc(zI-Tc)=I} and yields \code{norm(C Rf-Rc C) <= norm(Rc)*norm(defect)*norm(Rf)}; zero defect gives exact intertwining.
\par\smallskip\textbf{Boundary.} All inverse equations are supplied witnesses. Lean does not construct resolvents, prove spectral membership, or infer inverse existence from analytic hypotheses.
\medskip
\Needspace{7\baselineskip}
\subsection{Row 22: Smith phase inversion identity}
\label{row:22}
\textbf{Classification:} Stronger / scoped realization.
\par\smallskip\textbf{Declaration targets}\par
{\small\noindent\code{SD.smithPhase}\par}
{\small\noindent\code{SD.smithPhase_neg_one_eq_inv}\par}
\smallskip\textbf{Checked scope.} Over any field in which \code{2} is nonzero, the pulled-back phase formula at \code{omega = -1} equals \code{z^-1} for \code{z != 0}; this includes \code{Q}, \code{R}, and \code{C} and does not need the manuscript's additional \code{z != -1} condition.
\par\smallskip\textbf{Boundary.} Only the algebraic formula is checked. No Smith disk, Mobius-domain, radius-preservation, or physical impedance theorem is formalized.
\medskip
\Needspace{7\baselineskip}
\subsection{Row 23: Residue analysis and synthesis algorithm}
\label{row:23}
\textbf{Classification:} Specification.
\par\smallskip\textbf{Declaration targets}\par
{\small\noindent\code{PC.rangeSynthesis_rangeAnalysis}\par}
{\small\noindent\code{PC.rangeAnalysis_rangeSynthesis}\par}
{\small\noindent\code{PC.exactPacketCoordinates}\par}
{\small\noindent\code{NR.coordinateRefinementEquiv}\par}
\smallskip\textbf{Checked scope.} For finite complete projector families, range-restricted analysis and finite-sum synthesis are inverse linear maps; cross-level analysis/synthesis gives an exact coordinate equivalence.
\par\smallskip\textbf{Boundary.} This is the mathematical contract behind the array routine, not extracted interleaving code or a proof about Python storage shapes.
\medskip
\Needspace{7\baselineskip}
\subsection{Row 24: Weighted finite-precision inverse algorithm}
\label{row:24}
\textbf{Classification:} Specification / certificate-explicit.
\par\smallskip\textbf{Declaration targets}\par
{\small\noindent\code{WT.finiteGeometricInverse}\par}
{\small\noindent\code{WT.prunedWordInverse}\par}
{\small\noindent\code{WT.map_finiteGeometricInverse_eq_prunedWordInverse}\par}
{\small\noindent\code{WT.le_mul_ceilDiv}\par}
{\small\noindent\code{WT.prunedWordInverse_spec_of_minWeight}\par}
{\small\noindent\code{WT.prunedDiagonalTransferInverse}\par}
{\small\noindent\code{WT.finitePrecisionPrunedDiagonalTransferInverse_of_minWeight}\par}
{\small\noindent\code{WT.finitePrecisionPrunedDiagonalTransferInverse_of_weightedEdges}\par}
\smallskip\textbf{Checked scope.} The bounded geometric loop maps exactly to an explicit signed sum of retained noncommutative words. For \code{m>0}, Lean proves \code{L <= m * ceil(L/m)}, so the pseudocode choice \code{N=ceil(L/m)} meets the cutoff premise. Filtration data and \code{L<=m*N} prove that the pruned expression is a two-sided inverse; supplied \code{Dinv} and a normalized-transfer edge expansion lift it to the explicit pruned diagonal-plus-transfer inverse after applying \code{q}.
\par\smallskip\textbf{Boundary.} This is a mathematical specification. The manuscript calls the perturbation \code{H}; the Lean parameter name \code{K} has no fraction-field meaning. The result does not provide an extracted modular matrix inverse, basis representation, zero test, complexity bound, or proof of equivalence to the companion Python implementation.
\medskip
\Needspace{7\baselineskip}
\subsection{Row 25: Transfer-zeta walk construction algorithm}
\label{row:25}
\textbf{Classification:} Specification / automatic.
\par\smallskip\textbf{Declaration targets}\par
{\small\noindent\code{TZ.walkTuplesAtWeight}\par}
{\small\noindent\code{TZ.mem_walkTuplesAtWeight_iff}\par}
{\small\noindent\code{TZ.tuple_captured_iff}\par}
{\small\noindent\code{TZ.coefficientAlgorithm}\par}
{\small\noindent\code{TZ.coefficientAlgorithm_correct}\par}
{\small\noindent\code{TZ.exactWalkFirstEdgeEquiv}\par}
{\small\noindent\code{TZ.coefficientAlgorithm_firstEdge}\par}
{\small\noindent\code{TZ.pathZetaAlgorithm_trace}\par}
{\small\noindent\code{TZ.weightedAdjacency}\par}
{\small\noindent\code{TZ.weightedAdjacency_coeff}\par}
{\small\noindent\code{TZ.enumeratedPathZetaMatrix_eq_matrixPathZetaMatrix}\par}
{\small\noindent\code{TZ.pathZeta_eq_inverse_mul_transferNumerator}\par}
{\small\noindent\code{TZ.finite_graph_pathZeta_rationality}\par}
{\small\noindent\code{TZ.finite_graph_constant_coefficient_recurrence}\par}
\smallskip\textbf{Checked scope.} The coefficient loop is a finite \code{Finset} enumeration with exact capture and cardinality theorems. A direct trace wrapper pairs its semantic walk count with the first-edge recurrence. Finite matrix determinant/adjugate data are then constructed and proved equal to the enumerated series, yielding the normalized rational identity and eventual recurrence without a supplied rationality certificate.
\par\smallskip\textbf{Boundary.} The objects counted are directed walks, with repetitions allowed. The Lean polynomial matrix definitions are noncomputable, and no extracted determinant/adjugate program, cancellation or minimal-denominator routine, performance bound, or analytic convergence theorem is claimed.
\medskip
